\documentclass[../../../include/open-logic-section]{subfiles}

\begin{document}

\olfileid[fr]{sth}{z}{pairs}
\olsection{La paire}

Examinons maintenant l'axiome suivant :

\begin{axiom}[Paire]
Pour deux ensembles quelconques $a, b$, l'ensemble $\{a, b\}$ existe.
\[
	\forall a \forall b \exists P \forall x (x \in P \liff (x = a \lor x = b))
\]
\end{axiom}

Voici comment justifier cet axiome à partir de la conception
itérative. Supposons que $a$ soit disponible à l'étape $S$ et que
$b$ soit disponible à l'étape $T$. Soit $M$ la plus tardive des étapes
$S$ et $T$. Puisque $a$ et $b$ sont tous deux disponibles à l'étape
$M$, la collection $\{a,b\}$ peut être constituée en ensemble à
n'importe quelle étape postérieure à $M$.

Mais un instant !{} Pourquoi supposer qu'il \emph{existe} une étape
après $M$ ? S'il n'y en a aucune, notre justification échoue. Pour
justifier l'axiome de la paire, nous devons donc ajouter un principe
au récit donné dans la \olref[sth][z][story]{sec}, à savoir :
\begin{enumerate}
	\item[] \stagessucc. Il n'existe pas de dernière étape.
\end{enumerate}
Ce principe est-il justifié ? Rien, dans le récit de Shoenfield,
n'énonce \emph{explicitement} qu'il n'existe pas de dernière étape.
Même si, à strictement parler, il s'agit d'un ajout à notre récit,
ce principe s'accorde bien avec l'idée fondamentale selon laquelle
les ensembles se forment par étapes. Nous l'admettrons simplement
dans la suite, et admettrons donc aussi l'axiome de la paire.

Grâce à ce nouvel axiome, nous pouvons démontrer l'existence de bien
d'autres ensembles. Par exemple :

\begin{prop}\ollabel{prop:pairsconsequences}
Quels que soient les ensembles $a$ et $b$, les ensembles suivants existent :
	\begin{enumerate}
		\item\ollabel{singleton} $\{a\}$
		\item\ollabel{binunion} $a \cup b$
		\item\ollabel{tuples} $\tuple{a, b}$
	\end{enumerate}
\end{prop}

\begin{proof}
\olref{singleton}. Par l'axiome de la paire, $\{a, a\}$ existe ;
c'est $\{a\}$ par extensionnalité.

\olref{binunion}. Par l'axiome de la paire, $\{a, b\}$ existe.
Or $a \cup b = \bigcup \{a, b\}$ existe par l'axiome de l'union.

\olref{tuples}. D'après \olref{singleton}, $\{a\}$ existe. Par
l'axiome de la paire, $\{a, b\}$ existe. Enfin,
$\{\{a\}, \{a, b\}\} = \tuple{a, b}$ existe encore par l'axiome de la
paire.
\end{proof}

\begin{prob}
Montrez que, pour trois ensembles quelconques $a, b, c$, l'ensemble
$\{a, b, c\}$ existe.
\end{prob}

\begin{prob}
Montrez que, pour des ensembles quelconques $a_1, \ldots, a_n$, l'ensemble
$\{a_1, \ldots, a_n\}$ existe.
\end{prob}

%\begin{proof} Par deux applications de l'axiome de la paire,
%   $\{\{a_1, a_2\}, \{a_1, a_3\}\}$ existe. Par l'axiome de l'union
%   et l'extensionnalité,
%   $\bigcup \{\{a_1, a_2\}, \{a_1, a_3\}\} = \{a_1, a_2, a_3\}$ existe.
%   Répétez ce procédé autant de fois que nécessaire. \end{proof}

\end{document}
